% This file is part of the modular source for Stratified Relational Workspace Theory.
% !TeX root = ../Stratified_Relational_Workspace_Theory.tex
\theorypart{VIII}{Causal Efficacy and Reentry}

\topic{28}{Feasible Workspace Interventions}

A notation such as \(\operatorname{do}(W_t=w)\) is meaningful only if the
candidate workspace is an endogenous variable in a structural causal model or
if an intervention on underlying variables is shown to realize that value.
Abstractly setting a projection coordinate need not correspond to any possible
physical or computational manipulation.

\begin{definition}[Time-unrolled causal realization]\label{def:causal-realization}
A time-unrolled causal realization is a structural causal model
\[
\mathcal M
=
(\mathcal U,\mathcal V,F,\mathbb P_U)
\]
whose endogenous variables include the local states, messages, actions,
artifacts, target variables, and a measurable workspace extractor
\[
\pi_t:\mathsf S_t\to\mathsf W_t,
\qquad
W_t=\pi_t(S_t).
\]
The unrolled graph is acyclic even when the rolled process is recurrent.
\end{definition}

\begin{definition}[Feasible implementation intervention]
Let \(B_t\) denote declared background variables that should remain matched.
A family \(\{\iota_w:w\in\mathsf W_t^\circ\}\) is a feasible implementation of
workspace values on the experimental domain \(\mathsf W_t^\circ\) when:
\begin{enumerate}
\item each \(\iota_w\) acts only on declared manipulable parents, channels, or
 artifacts at or after time \(t\);
\item the intervened realization satisfies
\[
\pi_t(S_t^{\iota_w})=w
\qquad\text{almost surely};
\]
\item the background law is matched across compared interventions; and
\item the intervention does not directly modify \(Y_t^+\) except through
 paths included in the declared causal design.
\end{enumerate}
\end{definition}

Let \(D\) be a divergence or metric on probability laws. For a matched
background value \(b\), define
\[
\CE_{\Workspace}(w,w';b)
=
D\!\left(
\mathbb P^{\iota_w}(Y_t^+\in\cdot\mid B_t=b),
\mathbb P^{\iota_{w'}}(Y_t^+\in\cdot\mid B_t=b)
\right).
\]
For a declared comparison design \(\nu\) over \((w,w',b)\), define
\[
\boxed{
\eta_{\mathrm{causal}}^{\nu}
=
\mathbb E_{(w,w',b)\sim\nu}
\left[
\CE_{\Workspace}(w,w';b)
\right].
}
\]
The maximum
\[
\eta_{\mathrm{causal}}^{\max}
=
\sup_{w,w',b}\CE_{\Workspace}(w,w';b)
\]
may be useful for discovery, but a design-weighted effect is harder to inflate
by selecting a single exceptional pair after inspection.

\begin{proposition}[Coordinate invariance of intervention effects]
\label{prop:causal-gauge}
Let \(g:\mathsf W_t\to\mathsf W_t'\) be a measurable bijection and suppose the
reparametrized intervention family satisfies
\[
\iota_{g(w)}'=\iota_w.
\]
Then every divergence-based causal effect is preserved:
\[
\CE_{\Workspace'}(g(w),g(w');b)
=
\CE_{\Workspace}(w,w';b).
\]
\end{proposition}

\begin{proof}
The two labels designate the same underlying interventions and therefore induce
the same pair of target laws. Only the coordinate name has changed.
\end{proof}

This is the causal analogue of gauge invariance. It also shows why effects
should be attached to intervention implementations rather than to arbitrary
numerical coordinates.

\topic{29}{Causal Specificity and Transfer Mediation}

A causal effect on the target is not workspace-specific if an equal-cost
perturbation outside the candidate workspace produces the same effect. Let
\(c(\iota)\) be a declared intervention cost and let
\[
\mathcal I_{\Workspace}(r),
\qquad
\mathcal I_{\perp}(r)
\]
be matched families with cost \(r\), the first changing workspace coordinates
and the second changing complementary coordinates while preserving the
workspace value as closely as the implementation permits.

Define
\[
\eta_{\Workspace}(r)
=
\mathbb E_{\iota,\iota'\sim\mathcal I_{\Workspace}(r)}
D\!\left(
\mathbb P^{\iota}(Y_t^+),
\mathbb P^{\iota'}(Y_t^+)
\right),
\]
\[
\eta_{\perp}(r)
=
\mathbb E_{\jmath,\jmath'\sim\mathcal I_{\perp}(r)}
D\!\left(
\mathbb P^{\jmath}(Y_t^+),
\mathbb P^{\jmath'}(Y_t^+)
\right),
\]
and
\[
\boxed{
\Delta_{\mathrm{specific}}(r)
=
\eta_{\Workspace}(r)-\eta_{\perp}(r).
}
\]
A positive value supports causal concentration in the candidate workspace at
the declared perturbation scale.

If the complement is defined by \(\Gamma_\varepsilon(H)^\perp\), the result is
metric dependent. Under nonunitary gauge changes, either the metric must be
transported with the coordinates or interventions must be matched by a
coordinate-independent physical cost. Otherwise specificity can change merely
because the representation was rescaled.

\subtopic{Causal transfer nulls}

Let \(\jmath_E\) randomize selected communication edges while preserving their
marginal message laws, rates, and timing budget. The causal transfer gap is
\[
\boxed{
\Delta_{\mathrm{tr}}^{\operatorname{do}}(\jmath_E)
=
\Risk_{\ell,\mathbb P^{\jmath_E}}(V)
-
\Risk_{\ell,\mathbb P}(V).
}
\]
Unlike an observational null law, this quantity is anchored to an implemented
channel intervention. A positive effect is evidence that the intact alignment
of messages, rather than merely their marginal content, contributes to task
performance.

A stronger mediation test equalizes the workspace state after two distinct
channel interventions. Let \(\mu\) be a common imposed workspace distribution,
implemented by a feasible stochastic intervention in the sense of
Definition~\ref{def:causal-realization}. The notation below is shorthand for
that implementation. Define the residual mediation defect
\[
\delta_{\mathrm{med}}
=
\sup_{\jmath,\jmath'}
D\!\left(
\mathbb P^{\jmath,\operatorname{do}(W_t\sim\mu)}(Y_t^+),
\mathbb P^{\jmath',\operatorname{do}(W_t\sim\mu)}(Y_t^+)
\right).
\]
A small value means that once the workspace state is equalized, little of the
channel intervention's effect remains. This supports, but does not by itself
identify without the causal assumptions, the claim that \(W_t\) mediates the
relevant relational influence.

\topic{30}{Causal Recurrence, Loop-Mediated Reentry, and Relationality}

Temporal correlation is weaker than causal recurrence. Define the total
workspace recurrence effect at lag \(\tau\) by
\[
\boxed{
\eta_{\mathrm{rec}}(\tau)
=
\mathbb E_{(w,w',b)\sim\nu}
D\!\left(
\mathbb P^{\iota_w}(W_{t+\tau}\in\cdot\mid B_t=b),
\mathbb P^{\iota_{w'}}(W_{t+\tau}\in\cdot\mid B_t=b)
\right).
}
\]
A nonzero value shows that manipulating the present workspace changes a later
workspace state. It may, however, be explained by direct copying or a retained
memory register.

To isolate reentry, consider the unrolled causal graph and remove every direct
retention edge whose source and target are both workspace variables at
successive times. Let
\[
\mathcal M^{\setminus\mathrm{dir}}
\]
denote the resulting path-specific model, preserving paths that leave the
workspace through a nonworkspace participant, action, message, or artifact and
later return. Define
\[
\boxed{
\eta_{\mathrm{loop}}(\tau)
=
\mathbb E_{(w,w',b)\sim\nu}
D\!\left(
\mathbb P_{\mathcal M^{\setminus\mathrm{dir}}}^{\iota_w}
(W_{t+\tau}\in\cdot\mid B_t=b),
\mathbb P_{\mathcal M^{\setminus\mathrm{dir}}}^{\iota_{w'}}
(W_{t+\tau}\in\cdot\mid B_t=b)
\right).
}
\]
Then the distinctions are:
\[
\begin{aligned}
I(W_t;W_{t+\tau})>0
&\quad\text{temporal dependence only},\\
\eta_{\mathrm{rec}}(\tau)>0
&\quad\text{causal persistence},\\
\eta_{\mathrm{loop}}(\tau)>0
&\quad\text{causal exit-and-return reentry}.
\end{aligned}
\]
Path-specific effects require the usual consistency and no-unmeasured-confounding
assumptions appropriate to the implementation. In software systems, edge
interventions can often be literal. In human collaboration, these assumptions
are substantially harder to justify.

\begin{theorem}[Operational relationality, \(T_{13}\)]
\label{thm:operational-relationality}
Let \(\mathsf Y\) be finite, let \(P\) be the intact law at a fixed decision
time, and let \(W_t=\pi_t(\mathbf H_t)\) be a structurally supported candidate.
Let
\(\mathcal Q\subseteq\mathcal N_{\mathrm{pred}}(P)\) be a nonempty,
predeclared panel of validated relation-breaking nulls. Assume:
\begin{enumerate}
\item exact structural--operational alignment,
\[
I_P(Y_t^+;\mathbf H_t\mid W_t)=0;
\]
\item a positive node-level information gap,
\[
\delta_{\mathrm{node}}
=
\min_{i\in V}
I_P(Y_t^+;\mathbf H_t\mid\mathbf H_t^i)
>0;
\]
\item a positive robust matched-null gap,
\[
\delta_{\mathrm{null}}
=
\min_{Q\in\mathcal Q}
\left[
H_Q(Y_t^+\mid\mathbf H_t)
-
H_P(Y_t^+\mid\mathbf H_t)
\right]
>0;
\]
\item a feasible intervention design with
\[
\eta_{\mathrm{causal}}^{\nu}>0;
\]
\item positive matched causal specificity at some declared scale \(r\),
\[
\Delta_{\mathrm{specific}}(r)>0.
\]
\end{enumerate}
Then, relative to the declared decomposition, target, log loss, null panel, and
intervention family:
\begin{enumerate}
\item no statistic of any single-node history is sufficient for \(Y_t^+\), and
 \(k_0\geq2\);
\item \(W_t\) cannot be reconstructed almost surely from any single-node
 history;
\item every null \(Q\in\mathcal Q\) preserves all single-node predictive
 risks but has strictly worse full-system log risk than the intact law;
\item the candidate workspace is causally efficacious and its effect exceeds
 that of the declared matched off-workspace perturbations.
\end{enumerate}
If, in addition, \(\eta_{\mathrm{loop}}(\tau)>0\), the workspace is
loop-reentrant at lag \(\tau\).
\end{theorem}

\begin{proof}
By Theorem~\ref{thm:coalition-log-risk}, for each node \(i\),
\[
\Risk_{\log,P}(\{i\})-
\Risk_{\log,P}(V)
=
I_P(Y_t^+;\mathbf H_t\mid\mathbf H_t^i)
\geq
\delta_{\mathrm{node}}>0.
\]
Hence no single-node history, and therefore no statistic of one, is sufficient;
this gives \(k_0\geq2\).

Exact alignment and Theorem~\ref{thm:alignment-criterion} imply that the
canonical predictive state factors through \(W_t\). If
\(W_t=\psi_i(\mathbf H_t^i)\) almost surely for some node, then the canonical
state would also be a function of \(\mathbf H_t^i\), contradicting the positive
node gap through Proposition~\ref{prop:sufficiency-reconstruction}.

For every \(Q\in\mathcal Q\), Proposition~\ref{prop:local-risk-null} gives
\[
\Risk_{\log,Q}(\{i\})
=
\Risk_{\log,P}(\{i\})
\qquad\forall i.
\]
The assumed null gap gives
\[
\Risk_{\log,Q}(V)
>
\Risk_{\log,P}(V).
\]
Thus a description determined only by the collection of single-node predictive
marginals cannot reproduce the intact joint capacity: it would assign the same
local data to \(P\) and \(Q\) while their full-system risks differ.

Finally, the feasible intervention effect establishes causal efficacy relative
to the declared implementation, and the positive specificity gap shows that
the effect exceeds the matched complementary perturbation effect. The final
reentry statement is exactly the definition of \(\eta_{\mathrm{loop}}\).
\end{proof}

\begin{remark}[Scope of the theorem]
The conclusion is operational and indexed by
\[
(V,Y_t^+,\ell,\mathcal Q,\mathcal I,\nu,r).
\]
Changing the decomposition, target, loss, nulls, or intervention design may
change the verdict. The theorem does not establish metaphysical
irreducibility, autonomous agency, selfhood, or consciousness.
\end{remark}

\subtopic{Finite certification sequence}

A reproducible dynamic-causal analysis should proceed in the following order:
\begin{enumerate}
\item Freeze the node decomposition, accessible coalition histories, target,
 horizon, and proper scoring loss.
\item Estimate the structural projector on training data only and encode
 \(W_t\) without using held-out targets.
\item Cross-fit conditional target laws from the full history, the workspace,
 and every declared coalition. Report alignment and coalition risk gaps.
\item Construct the null panel before inspecting null performance. Verify
 local-predictive matching and the intended destruction of transfer
 organization separately.
\item Implement feasible workspace and complementary interventions with a
 common background and cost design. Report both causal effect and
 specificity.
\item Test recurrence and then cut direct retention paths to test loop-mediated
 reentry.
\item Report confidence intervals or finite-sample uncertainty for every gate;
 do not substitute training accuracy for Bayes-risk or causal estimates.
\end{enumerate}

The resulting certification report is
\[
\boxed{
\begin{aligned}
\mathcal C_{\mathrm{dyn}}(W_t)
=\bigl(&
\delta_{\mathrm{align}}^{\mathcal I},
\{\Delta_A^\ell\}_{A\subseteq V},
 k_\varepsilon,
\boldsymbol\Delta_{\mathrm{tr}}(\mathcal Q),\\
&\eta_{\mathrm{causal}}^\nu,
\Delta_{\mathrm{specific}}(r),
\eta_{\mathrm{rec}}(\tau),
\eta_{\mathrm{loop}}(\tau)
\bigr).
\end{aligned}
}
\]
No single coordinate certifies a relational workspace.
